\documentclass[12pt]{article}

\input{../../../_assets/preambulo}

\begin{document}

    % 1. Foto de fondo
    % 2. Título
    % 3. Encabezado Izquierdo
    % 4. Color de fondo
    % 5. Coord x del titulo
    % 6. Coord y del titulo
    % 7. Fecha

    
	\input{../../../_assets/portada}
	\portadaExamenFotoDif{../../sapienza.jpg}{Computer Vision\\Examen IV}{Computer Vision. Examen IV}{MidnightBlue}{-8}{28}{2025-2026}{Joaquín Avilés de la Fuente}

    \begin{description}
        \item[Asignatura] Computer Vision.
        \item[Curso Académico] 2025-2026.
        \item[Grado] Grado en Informática.
        \item[Grupo] Grupo A.
        \item[Profesor] Irene Amerini.
        \item[Descripción] Examen final escrito de Computer Vision.
        \item[Fecha] junio de 2026.
        \item[Duración] 60 min.
    
    \end{description}
    \newpage

    % ------------------------------------

    \begin{observacion}
        The total sum of all the exercises is 32 points, although the maximum mark of the exam is 30.
    \end{observacion}

    \begin{ejercicio}[8 points]
    Discuss the role of image derivatives in edge and corner detection. In your answer, explain:
    \begin{enumerate}
        \item[a)] The meaning of first-order and second-order image derivatives.
        \item[b)] How gradient magnitude and gradient direction are used for edge detection.
        \item[c)] How second-order derivatives can be used to identify intensity transitions.
        \item[d)] Why are derivative-based methods sensitive to noise and how can mitigate that.
    \end{enumerate}
    \end{ejercicio}

    \begin{ejercicio}[8 points]
        You are given a set of images of the same building taken from different viewpoints, some images contain repeated structures such as windows, columns, or similar textures. You need to build a robust image matching pipeline for 3D reconstruction. Describe:
        \begin{enumerate}
            \item[a)] Which local features and descriptors you would use and why.
            \item[b)] How you would match descriptors between image pairs.
            \item[c)] How you would remove incorrect matches caused by repeated patterns.
        \end{enumerate}
    \end{ejercicio}

    \begin{ejercicio}[8 points]
        Explain the concept of homography and its use in image alignment. In your answer, discuss:
        \begin{enumerate}
            \item[a)] What a homography matrix represents geometrically.
            \item[b)] Under which assumptions a homography can correctly map points between two images.
            \item[c)] How corresponding feature points can be used to estimate a homography.
        \end{enumerate}
    \end{ejercicio}

    \begin{ejercicio}[8 points]
        Describe the process of recovering 3D structure from two calibrated images. In your answer, include:
        \begin{enumerate}
            \item[a)] The role of camera calibration and intrinsic/extrinsic parameters.
            \item[b)] How triangulation is used to estimate the 3D position of scene points.
        \end{enumerate}
    \end{ejercicio}
\end{document}